\documentclass[12pt, fleqn, openany]{studyGuide}
\setstretch{1.4}
\enableInlineReview
\renewcommand{\VMBookTitle}{Audit Review}
\renewcommand{\VMBookSubtitle}{Grade 7 --- Ch Graphic Review --- Changed Questions}
\renewcommand{\VMFooterURL}{https://viewmath.com/grade7}
\renewcommand{\VMFooterURLDisplay}{ViewMath.com/Grade7}
\begin{document}

\begin{center}
\begin{tcolorbox}[
	enhanced,
	colback=white,
	colframe=funTeal!60,
	boxrule=1.5pt,
	arc=10pt,
	outer arc=10pt,
	width=0.9\linewidth,
	halign=center,
	top=5mm, bottom=5mm,
]
	{\Large\bfseries\sffamily\color{funTealDark}Audit Review}\\[2mm]
	{\large\sffamily\color{funGrayDark}Grade 7 --- Ch Graphic Review --- Changed Questions}\\[2mm]
	{\normalsize\sffamily\color{funGrayDark}Verify corrections below}
\end{tcolorbox}
\end{center}

\newpage
\section*{9-5-q26 \hfill \normalsize\texttt{tests\_questions\_bank\_2/topics/ch09-05-sample-spaces-for-compound-events.tex}}
\begin{practiceQuestion}{9-5-q26}{gsa}
\begin{questionText}
Draw the sample space as an organized list for the compound event: flip a coin and pick a card from $\{1, 2, 3\}$.

\begin{center}
\begin{tikzpicture}[x=1cm, y=1cm, >=Latex, font=\small]
  \node at (0,4.4) {Start};
  \node at (-2.3,2.8) {H};
  \node at (2.3,2.8) {T};
  \node at (-3.9,0.9) {H1};
  \node at (-2.3,0.9) {H2};
  \node at (-0.7,0.9) {H3};
  \node at (0.7,0.9) {T1};
  \node at (2.3,0.9) {T2};
  \node at (3.9,0.9) {T3};

  \draw[thick,->] (0,4.12) -- (-2.3,3.06);
  \draw[thick,->] (0,4.12) -- (2.3,3.06);

  \draw[thick,->] (-2.3,2.48) -- (-3.9,1.18);
  \draw[thick,->] (-2.3,2.48) -- (-2.3,1.18);
  \draw[thick,->] (-2.3,2.48) -- (-0.7,1.18);
  \draw[thick,->] (2.3,2.48) -- (0.7,1.18);
  \draw[thick,->] (2.3,2.48) -- (2.3,1.18);
  \draw[thick,->] (2.3,2.48) -- (3.9,1.18);

  \node[font=\scriptsize] at (-3.15,1.95) {1};
  \node[font=\scriptsize] at (-2.3,1.9) {2};
  \node[font=\scriptsize] at (-1.45,1.95) {3};
  \node[font=\scriptsize] at (1.45,1.95) {1};
  \node[font=\scriptsize] at (2.3,1.9) {2};
  \node[font=\scriptsize] at (3.15,1.95) {3};
\end{tikzpicture}
\end{center}

How many outcomes have tails and an even-numbered card?
\end{questionText}
\correctAnswer{$1$}
\explanation{From the tree diagram, the outcomes with tails are T1, T2, T3. The only even number is $2$, so the only outcome with tails and an even card is T2. That gives $1$ outcome.}
\end{practiceQuestion}

\newpage
\section*{9-6-q23 \hfill \normalsize\texttt{tests\_questions\_bank\_2/topics/ch09-06-finding-probabilities-of-compound-events.tex}}
\begin{practiceQuestion}{9-6-q23}{gmc}
\begin{questionText}
The tree diagram below shows the outcomes for choosing a sandwich type and a side dish.

\begin{center}
\begin{tikzpicture}[x=1cm, y=1cm, >=Latex, font=\small]
  \node at (0,4.4) {Start};
  \node at (-3.2,2.8) {Ham};
  \node at (0,2.8) {Turkey};
  \node at (3.2,2.8) {Veggie};
  \node at (-4.5,0.9) {Chips};
  \node at (-1.9,0.9) {Fruit};
  \node at (-1.3,0.9) {Chips};
  \node at (1.3,0.9) {Fruit};
  \node at (1.9,0.9) {Chips};
  \node at (4.5,0.9) {Fruit};

  \draw[thick,->] (0,4.12) -- (-3.2,3.08);
  \draw[thick,->] (0,4.12) -- (0,3.08);
  \draw[thick,->] (0,4.12) -- (3.2,3.08);

  \draw[thick,->] (-3.2,2.48) -- (-4.5,1.18);
  \draw[thick,->] (-3.2,2.48) -- (-1.9,1.18);
  \draw[thick,->] (0,2.48) -- (-1.3,1.18);
  \draw[thick,->] (0,2.48) -- (1.3,1.18);
  \draw[thick,->] (3.2,2.48) -- (1.9,1.18);
  \draw[thick,->] (3.2,2.48) -- (4.5,1.18);
\end{tikzpicture}
\end{center}

If each combination is equally likely, what is the probability of getting Turkey with Fruit?
\end{questionText}
\choiceA{$\dfrac{1}{2}$}
\choiceB{$\dfrac{1}{3}$}
\choiceC{$\dfrac{1}{5}$}
\choiceD{$\dfrac{1}{6}$}
\correctAnswer{D}
\explanation{The tree diagram shows $3 \times 2 = 6$ equally likely outcomes. ``Turkey with Fruit'' is $1$ outcome, so $P = \frac{1}{6}$.}
\end{practiceQuestion}

\newpage
\section*{9-6-q26 \hfill \normalsize\texttt{tests\_questions\_bank\_2/topics/ch09-06-finding-probabilities-of-compound-events.tex}}
\begin{practiceQuestion}{9-6-q26}{gsa}
\begin{questionText}
The tree diagram shows outcomes for choosing an appetizer and an entree at a restaurant.

\begin{center}
\begin{tikzpicture}[x=0.9cm, y=0.9cm, >=Latex, font=\small]
  \node at (0,5.1) {Start};
  \node at (-2.8,3.4) {Soup};
  \node at (2.8,3.4) {Bread};
  \node at (-5.8,1.1) {Steak};
  \node at (-4.0,1.1) {Pasta};
  \node at (-2.2,1.1) {Fish};
  \node at (-0.4,1.1) {Salad};
  \node at (1.4,1.1) {Steak};
  \node at (3.2,1.1) {Pasta};
  \node at (5.0,1.1) {Fish};
  \node at (6.8,1.1) {Salad};

  \draw[thick,->] (0,4.78) -- (-2.8,3.66);
  \draw[thick,->] (0,4.78) -- (2.8,3.66);

  \draw[thick,->] (-2.8,3.08) -- (-5.8,1.38);
  \draw[thick,->] (-2.8,3.08) -- (-4.0,1.38);
  \draw[thick,->] (-2.8,3.08) -- (-2.2,1.38);
  \draw[thick,->] (-2.8,3.08) -- (-0.4,1.38);
  \draw[thick,->] (2.8,3.08) -- (1.4,1.38);
  \draw[thick,->] (2.8,3.08) -- (3.2,1.38);
  \draw[thick,->] (2.8,3.08) -- (5.0,1.38);
  \draw[thick,->] (2.8,3.08) -- (6.8,1.38);
\end{tikzpicture}
\end{center}

If each combination is equally likely, what is the probability of choosing Bread as the appetizer and NOT choosing Steak? Express your answer as a simplified fraction.
\end{questionText}
\correctAnswer{$\frac{3}{8}$}
\explanation{Total outcomes: $2 \times 4 = 8$. Bread with something other than Steak: Bread-Pasta, Bread-Fish, Bread-Salad — $3$ outcomes. $\frac{3}{8}$.}
\end{practiceQuestion}

\end{document}